\section{Period dynamics and the conclusion}
\label{dyn:section}

We now pass from the two open regions constructed above to every
K3 surface. The distinction between no rational direction, one
rational direction, and a rational plane is essential. In
particular, density of an Oka locus alone would not justify this
step. The corrected period-orbit analysis is due to Verbitsky
\cite{VerbitskyErratum}; the recent dense-period application in
\cite{XieZhao} uses the same broad strategy of geometric flexibility, arithmetic
period dynamics, and Torelli theory. Here the two preceding geometric
constructions supply open regions for both nonrational orbit types.
We give the particular group argument needed for this transfer.

Write $V=\Lambda\otimes\R$, and let
$G=\operatorname{SO}_0(V)$ be the identity component. Passing to
finite-index subgroups, the integral isometry group gives an
arithmetic lattice $\Gamma\subset G$ by the arithmetic lattice
theorem \cite{BorelHarishChandra}.

\begin{lemma}\label{dyn:orbits}
Let $P\in\mathcal D$.
\begin{enumerate}
\item If $P\cap(\Lambda\otimes\Q)=\{0\}$, its integral-isometry
orbit is dense in $\mathcal D$.
\item If $P\cap(\Lambda\otimes\Q)=\Q v$ for a primitive integral
vector $v$, its orbit under integral isometries fixing $v$ is
dense in $\mathcal D_v$.
\end{enumerate}
\end{lemma}

\begin{proof}
We first work with an ordered orthonormal positive frame of $P$.
Its connected pointwise stabilizer is
$H=\operatorname{SO}_0(1,19)$. This group is generated by
one-parameter unipotent subgroups. Ratner's orbit closure
theorem \cite{Ratner}, applied to $H\cdot e\Gamma$ in
$G/\Gamma$, gives
\[
 \overline{H\cdot e\Gamma}=S\cdot e\Gamma
\]
for a connected closed subgroup $H\subset S\subset G$ such
that $S\cap\Gamma$ is a lattice in $S$. Thus $H$ is the
pointwise stabilizer of the actual plane $P$ throughout, and the lattice
matrices in every intermediate group are rational in the
original lattice coordinates.

Put $W=P^\perp$. As a module for
$\mathfrak h=\mathfrak{so}(W)$, the orthogonal block-matrix
decomposition is
\[
 \mathfrak{so}(V)
 =\mathfrak h\oplus\mathfrak{so}(P)\oplus(W\otimes P).
\]
The standard real representation $W$ is absolutely irreducible.
Thus an intermediate Lie algebra obtains its mixed part by
choosing a subspace $A\subset P$, with mixed module $W\otimes A$.
If $\dim A=2$, brackets give the full Lie algebra. If
$\dim A=1$, they give
$\mathfrak{so}(W\oplus A)$, the stabilizer of the positive line
$A^\perp\cap P$. Including the nonzero rotation algebra of $P$
in this case generates the missing mixed summand and hence the
full algebra. If $A=0$, the only possibilities are
$\mathfrak h$ and
$\mathfrak h\oplus\mathfrak{so}(P)$.

The proper alternatives force rational directions in $P$.
For the stabilizer of one positive line, Borel density
\cite{BorelDensity} makes its lattice Zariski dense in that
stabilizer. As the lattice consists of rational matrices, its
Zariski closure is defined over $\Q$. Its fixed line is therefore
rational. The same reasoning for $H$ makes its fixed plane
rational.

There is also no difficulty from the compact rotation factor in
the remaining algebra. Up to finite isogeny the corresponding
connected group is
$\operatorname{SO}_0(W)\times\operatorname{SO}(P)$.
Projection to the first factor is proper because its kernel is
compact. Thus the image of the lattice meets each compact set
in only finitely many points and is discrete. The finite
invariant measure on the original quotient pushes forward to
a finite invariant measure on the projected quotient, so the
image is a lattice there. Borel density implies that the
Lie algebra of the Zariski closure of the original lattice
surjects onto $\mathfrak h$. Its derived algebra is exactly
$\mathfrak h$: it is contained there because the other factor
is commutative, and it maps onto $\mathfrak h$ because that
algebra is simple. The Zariski closure and its derived algebra
are defined over $\Q$, so the fixed space of the derived
algebra, namely $P$, is rational. Here the rationality of the
Zariski closure of rational matrices follows directly from the
vanishing equations: in each bounded degree their linear
conditions have rational coefficients and hence a rational
basis of solutions.

If $P$ has no rational direction, none of these proper
alternatives is possible. Ratner's closure is consequently the
whole homogeneous space. Equivalently $H\Gamma$ is dense in
$G$. Inversion makes $\Gamma H$ dense as well, and projection
to $G/H$ makes the $\Gamma$-orbit of the chosen frame dense in
frame space.
The projection from frame space to $\mathcal D$ proves the first
assertion. The use of the frame stabilizer $H$, rather than
$H\times\operatorname{SO}(2)$, is what permits the application
of Ratner's theorem.

For the second assertion, restrict to $v^\perp$, of signature
$(2,19)$. The integral stabilizer of $v$ is arithmetic and gives
a lattice in $G_v=\operatorname{SO}_0(v^\perp)$. More explicitly,
the subgroup of integral isometries of the orthogonal lattice
acting trivially on its discriminant group extends to the full
K3 lattice while fixing $v$ \cite[Chapter~14, Section~2]{Huybrechts},
so the integral stabilizer has
finite index in the relevant orthogonal arithmetic group.

An oriented plane containing $v$ is equivalent to a positive
unit vector $u\in v^\perp$: choose the sign of $u$ using the
given plane orientation. The connected stabilizer of $u$ is
again $H=\operatorname{SO}_0(1,19)$. This time
\[
 \mathfrak{so}(2,19)=\mathfrak h\oplus W
\]
as an $\mathfrak h$-module, so irreducibility shows that $H$ is
maximal at Lie algebra level. Ratner's alternatives are a closed
$H$-orbit or a dense orbit. In the closed case Borel density
makes the line $\R u$ rational in the rational quadratic space
$v^\perp$. This would give a second rational direction in $P$,
contrary to the hypothesis. The orbit is therefore dense.
The connected group $G_v$ is transitive on positive unit
vectors, so the resulting density is throughout $\mathcal D_v$.
\end{proof}

\begin{proof}[Proof of Theorem~\ref{thm:main}]
Choose a marking of $X$ and let $P\in\mathcal D$ be its period.
The rational dimension of $P\cap(\Lambda\otimes\Q)$ is zero,
one, or two. We show in each case that $X$ has $L$ everywhere.

If it is zero, Lemma~\ref{dyn:orbits} makes the marked period
orbit dense in $\mathcal D$. It meets the nonempty open subset
provided by Proposition~\ref{deg:open}. Re-marking does not change
the complex surface. Moreover K3 surfaces with the same marked
period are isomorphic after their markings are forgotten: signs
and reflections in algebraic roots match their K\"ahler
chambers, and global Torelli then supplies an isomorphism.
Thus the all-point property $L$ in that open region holds on $X$.

If the rational dimension is one, let $v$ be a primitive
integral generator of that direction. It is positive because
$P$ is positive. The second part of Lemma~\ref{dyn:orbits}
allows us to re-mark within the stabilizer of $v$ so that the
period lies in the relatively open set $\mathcal O_v$ of
Theorem~\ref{rat:open}. The same Torelli argument gives $L$
everywhere on $X$. Notice that this step uses an open subset
inside the fixed-v locus itself.

If the rational dimension is two, $P$ is a rational plane.
Its rational orthogonal complement has signature $(1,19)$.
The positive cone is open and nonempty, so it contains a
rational vector, which can be multiplied by an integer to give
a positive integral $(1,1)$-class. The K3 projectivity criterion
implies that $X$ is projective. Proposition~\ref{prj:strips}
therefore gives $L$ at every point of $X$.

Theorem~\ref{ana:cap} now applies to $X$, proving the stated
approximation statement with all its quantifiers. The estimates
there are compact-uniform estimates for maps into $X$.
On the compact target, any two smooth Hermitian distances are
uniformly equivalent, so the conclusion holds for the
arbitrary metric specified in the statement.
\end{proof}
